% =====================================================================
%  1bat-ud3-5.tex  —  HORA 5: Primer contacte amb GeoGebra: famílies de funcions
%  (sense preàmbul: s'inclou des de main.tex)
% =====================================================================
\horatitol{Sessió 5 — Primer contacte amb GeoGebra: famílies de funcions}%
{Aprendre les accions bàsiques de GeoGebra i explorar com el pendent $m$ i el terme independent $n$ transformen una recta.}

\subsection{Idea clau}

A la sessió anterior vam veure que tota recta s'escriu $y=mx+n$ i que $m$ i $n$ tenen
un sentit social (ritme i punt de partida). Avui ho comprovem \textbf{en moviment}
amb \textbf{GeoGebra}: en lloc de dibuixar a mà, deixarem que l'ordinador dibuixi i
\emph{nosaltres} mourem $m$ i $n$ per veure com reacciona la gràfica. Aquesta sessió
és pràctica: el full guia us porta pas a pas. No cal saber-ne abans.

\begin{idea}
\textbf{Tres accions bàsiques de GeoGebra}
\begin{enumerate}[leftmargin=1.6em, itemsep=2pt]
  \item \textbf{Introduir una funció:} a la \emph{barra d'entrada} escriviu, per
        exemple, \texttt{f(x)=2x+1} i premeu \textsf{Retorn}. La recta apareix sola.
  \item \textbf{Moure i ajustar la finestra:} arrossegueu el full amb el ratolí per
        desplaçar-vos; useu la roda (o els botons $+$/$-$) per acostar o allunyar.
  \item \textbf{Crear un \emph{lliscador}} (paràmetre que es pot moure): escriviu
        \texttt{m=1} i premeu \textsf{Retorn}; apareix un lliscador. Després escriviu
        \texttt{f(x)=m*x+1} i, en moure el lliscador \texttt{m}, la recta canviarà
        \textbf{en directe}.
\end{enumerate}
\end{idea}

\subsection{Exemples}

\begin{exemple}[variar el pendent $m$ (mateixa $n$)]
Introduïu a GeoGebra la família $f(x)=mx+2$ amb un lliscador per a \texttt{m}. En
moure'l, totes les rectes passen pel mateix punt $(0,2)$, però s'\textbf{inclinen}
de manera diferent. Això és el que veureu (tres valors de $m$ amb $n=2$ fix):

\begin{center}
\begin{tikzpicture}
\begin{axis}[udaxis, width=9cm, height=5.4cm,
    xlabel={\small $x$}, ylabel={\small $y$},
    xmin=-4.5, xmax=4.5, ymin=-3, ymax=8,
    xtick={-4,-2,0,2,4}, ytick={-2,0,2,4,6}]
  \addplot[udblue, domain=-4:4, samples=2]{0.5*x+2};
  \addplot[udnavy, domain=-4:3, samples=2]{2*x+2};
  \addplot[udgrey, domain=-4:4, samples=2]{-1*x+2};
  \addplot[black, only marks, mark=*, mark size=1.5pt] coordinates {(0,2)};
  \node[font=\scriptsize, anchor=south west] at (axis cs:0,2) {$(0,2)$};
  \node[udnavy, font=\scriptsize, anchor=west] at (axis cs:2.3,6.6) {$m=2$};
  \node[udblue, font=\scriptsize, anchor=west] at (axis cs:3.3,3.7) {$m=0,5$};
  \node[udgrey, font=\scriptsize, anchor=east] at (axis cs:-3.2,5.2) {$m=-1$};
\end{axis}
\end{tikzpicture}

\figcap{La família $f(x)=mx+2$: canviant $m$, la recta s'inclina, però sempre passa per $(0,2)$.}
\end{center}

\noindent Com més gran és $m$, més \textbf{dreta} (inclinada) puja la recta; si $m$
és negatiu, \textbf{baixa}.
\end{exemple}

\begin{exemple}[variar el terme independent $n$ (mateixa $m$)]
Ara fixeu el pendent i moveu \texttt{n}: introduïu $f(x)=x+n$ amb un lliscador per a
\texttt{n}. Totes les rectes tenen la \textbf{mateixa inclinació} (són paral·leles),
però estan \textbf{a altures diferents}: pugen o baixen segons $n$.

\begin{center}
\begin{tikzpicture}
\begin{axis}[udaxis, width=9cm, height=5cm,
    xlabel={\small $x$}, ylabel={\small $y$},
    xmin=-4.5, xmax=4.5, ymin=-4, ymax=7,
    xtick={-4,-2,0,2,4}, ytick={-3,0,3}]
  \addplot[udblue, domain=-4:4, samples=2]{x+3};
  \addplot[udnavy, domain=-4:4, samples=2]{x+0};
  \addplot[udgrey, domain=-4:4, samples=2]{x-3};
  \node[udblue, font=\scriptsize, anchor=south] at (axis cs:1.6,5.2) {$n=3$};
  \node[udnavy, font=\scriptsize, anchor=south] at (axis cs:2.6,3.2) {$n=0$};
  \node[udgrey, font=\scriptsize, anchor=north] at (axis cs:3.4,-0.2) {$n=-3$};
\end{axis}
\end{tikzpicture}

\figcap{La família $f(x)=x+n$: canviant $n$, la recta puja o baixa sense canviar la inclinació.}
\end{center}
\end{exemple}

\subsection{Exercicis resolts}

\begin{exresolt}[anticipar el moviment de la recta]
A GeoGebra teniu $f(x)=mx+2$ amb \texttt{m} en un lliscador. Sense tocar res, digeu
què passarà amb la recta si moveu \texttt{m} de $1$ a $3$, i comproveu-ho després.
\solucio
Augmentar \texttt{m} fa créixer el \textbf{ritme de pujada}: la recta es torna més
\textbf{inclinada} (puja més de pressa). El punt $(0,2)$ no es mou, perquè $n=2$
segueix igual; només canvia la inclinació. En passar de $m=1$ a $m=3$, la recta
gira al voltant de $(0,2)$ fent-se més vertical.
\resposta{La recta s'inclina més (puja més ràpid) i continua passant per $(0,2)$.}
\end{exresolt}

\begin{exresolt}[traduir un canvi social a un paràmetre]
Una tarifa social és $C(x)=0,10x+5$ (€), on $x$ són els minuts. Si l'administració
decideix \textbf{apujar la part fixa} de $5\eur$ a $8\eur$ sense tocar el preu per
minut, quin paràmetre canvia a GeoGebra i com es mou la gràfica?
\solucio
La «part fixa» és el terme independent $n$. Passa de $n=5$ a $n=8$: canvia
\textbf{només $n$}, no $m$. A la gràfica, la recta \textbf{puja} $3$ unitats sencera,
sense canviar la inclinació (queda paral·lela a l'anterior). El preu per minut
($m=0,10$) no s'altera.
\resposta{Canvia $n$ (de $5$ a $8$); la recta puja $3$ i es manté paral·lela.}
\end{exresolt}

\subsection{Exercicis per fer}

\noindent\textit{Fes aquestes exploracions a GeoGebra i registra el que observes a la
teva llibreta o al full guia.}

\begin{exercicis}
\begin{exlist}
  \item \textbf{(Exploració)} Introdueix $f(x)=mx+1$ amb un lliscador per a \texttt{m}.
        Mou \texttt{m} i completa: quan $m>0$ la recta \rule{2.5cm}{0.4pt}; quan
        $m<0$ la recta \rule{2.5cm}{0.4pt}; quan $m=0$ la recta és \rule{2.5cm}{0.4pt}.
  \item \textbf{(Exploració)} Introdueix $f(x)=2x+n$ amb un lliscador per a \texttt{n}.
        Mou \texttt{n} entre $-3$ i $3$ i descriu amb les teves paraules què li passa
        a la recta. On talla l'eix vertical en cada cas?
  \item Anticipa \emph{abans} de fer-ho a GeoGebra: si tens $f(x)=3x+1$ i la canvies
        per $g(x)=3x+4$, com es mourà la gràfica? I si la canvies per $h(x)=5x+1$?
        Comprova-ho després.
  \item Escriu a GeoGebra aquestes dues funcions alhora i digues quina puja més de
        pressa i quina té el punt de partida més alt:
        \[
        f(x)=2x+1 \qquad\text{i}\qquad g(x)=x+4.
        \]
  \item \textbf{(Interpretació social)} Un impost es modelitza amb $T(x)=0,2x+10$.
        \begin{enumerate}[label=\alph*), leftmargin=1.6em, topsep=2pt]
          \item Si el govern n'\textbf{apuja el pendent} a $0,3$, què vol dir
                socialment? Com es mou la gràfica?
          \item Si en lloc d'això \textbf{apuja la part fixa} a $20$, què canvia per a
                qui té ingressos baixos? Com es mou la gràfica?
        \end{enumerate}
  \item \textbf{(Síntesi)} Escriu, amb les teves paraules, una frase que resumeixi
        \emph{què fa cada paràmetre}: una per a $m$ i una per a $n$. Aquesta serà la
        conclusió que tancarem en comú a classe.
\end{exlist}
\end{exercicis}

% ---------------------------------------------------------------------
\fullsolucions{Sessió 5}
\begin{solucions}
\begin{sollist}
  \item Quan $m>0$ la recta \textbf{creix} (puja cap a la dreta); quan $m<0$
        \textbf{decreix} (baixa); quan $m=0$ és \textbf{horitzontal}.
  \item En moure \texttt{n}, totes les rectes són paral·leles (mateix pendent $2$)
        però pugen o baixen; cadascuna talla l'eix vertical al punt $(0,n)$.
  \item De $f(x)=3x+1$ a $g(x)=3x+4$: la recta \textbf{puja} $3$ unitats, paral·lela
        (només canvia $n$). A $h(x)=5x+1$: mateix punt de partida $(0,1)$ però
        \textbf{més inclinada} (puja més de pressa, canvia $m$).
  \item $f(x)=2x+1$ puja més de pressa (pendent $2>1$); $g(x)=x+4$ té el punt de
        partida més alt ($n=4>1$).
  \item (a) Apujar el pendent significa que l'impost creix més de pressa amb els
        ingressos (els qui més tenen paguen proporcionalment més); la recta s'inclina
        més. (b) Apujar la part fixa fa que \emph{tothom}, també qui té ingressos
        baixos, pagui més des del principi; la recta puja sencera (paral·lela).
  \item Resposta oberta. Idea esperada: «$m$ controla la \emph{inclinació} (com de
        ràpid puja o baixa la recta)» i «$n$ controla l'\emph{altura} (on talla l'eix
        vertical, el punt de partida)».
\end{sollist}
\end{solucions}
